\section{案例回放：Prime 定理的多策略搜索}

\subsection{问题拆解与语义化}
Hubert 把一个 prime 定理拆解成三个子 goal：1）从定义出发构造候选数；2）用 lemma 联结 prime 与 modulus；3）用 tactic 列出推理 steps。每个子 goal 都在 Lean 中形成独立的 goal tree，AlphaProof 会先在 Mathlib 中查找相似 lemma，再调用 policy 生成 tactic。

\begin{knowledgebox}{Prime 定理的语义纹理}
1）定义层：prime number 需要从 basic arithmetic 过渡到 advanced number theory；2）策略层：不同 lemma 如 \texttt{mod\_pow}、\texttt{order\_dvd} 分别用不同 tactic；3）治理层：每个 unsuccessful tactic 都写入 trace，供 human reviewer 复审。
\end{knowledgebox}

\subsection{策略组合与 Chunk gating}
在实际执行中，AlphaProof 同时维持三条 search track：heuristic-run （低成本），beam-run（高保真），MCTS-run（深度 polish）。chunk gating 表里记录了每条 track 的 token budget、overlap、entropy drift；当 heuristic-run 勾勒出可能解时，beam-run 会接手；若 drift 超过阈值则降级到 MCTS-run，确保 eventual proof 的正确性。

\begin{table}[H]
\centering
\begin{tabular}{p{3.5cm}p{3.5cm}p{5cm}}
\toprule
Track & Budget & 说明 \\
\midrule
Heuristic-run & 128 tokens & 快速探索，生成 proof sketch；当 heatmap 低于 0.4 后让 beam 接入 \\
Beam-run & 256 tokens & 保守扩展 top-k candidate，加入 lemmas \\
MCTS-run & 512 tokens & 深度 polish，triggered by drift gate \\
\bottomrule
\end{tabular}
\caption{Prime 定理 search track 与 gating}
\end{table}

\begin{warningbox}{实时调度的注意事项}
Chunk gating 不能只是 threshold table，还要根据 proof term 长度调整 budget；如果 drift gate 频繁触发，说明 Beam track parameter 过 aggressive。
\end{warningbox}

\subsection{本章小结}
Prime 定理案例展现了 multi-track search 与 chunk gating gate 的配合，让 RL policy 在数学定理上既能快速探索又不失准确性；Trace log 让 drift gate 遇到的失败成为治理的素材。
\newpage

